\section{Salience Detection}

Salience quantifies the degree to which an incoming semantic state is novel,
unexpected, or behaviorally relevant. In the MEM8-Geometric-Architecture,
salience is defined as a geometric signal combining statistical novelty and
temporal change. This dual formulation ensures that both out-of-distribution
events and rapid semantic shifts are captured.

Let $\mathbf{x}_t \in \mathbb{R}^d$ denote the fused semantic vector at time $t$.
We define two complementary components of salience: novelty and temporal change.

\subsection{Statistical Novelty}

Novelty is measured using the Mahalanobis distance, which evaluates how far
$\mathbf{x}_t$ deviates from the expected distribution of semantic states.

Let $\bm{\mu}$ be the mean semantic vector and $\Sigma$ the covariance matrix.
Then the novelty score is:

\[
N(\mathbf{x}_t) =
\sqrt{(\mathbf{x}_t - \bm{\mu})^\top \Sigma^{-1}
(\mathbf{x}_t - \bm{\mu})}
\]

This metric provides a geometry-aware measure of deviation, sensitive to
directions of high and low variance in the semantic manifold.

\subsection{Temporal Change}

Temporal change captures the magnitude of semantic drift between consecutive
states:

\[
T(\mathbf{x}_t) = \|\mathbf{x}_t - \mathbf{x}_{t-1}\|
\]

This term detects rapid transitions, even when the new state is statistically
common. It is essential for identifying abrupt shifts in modality, context, or
intent.

\subsection{Combined Salience}

The full salience score is a weighted combination of novelty and temporal
change:

\[
S(\mathbf{x}_t) =
\alpha\, N(\mathbf{x}_t) + (1 - \alpha)\, T(\mathbf{x}_t)
\]

where $\alpha \in [0,1]$ controls the balance between distributional novelty
and temporal dynamics.

\subsection{Geometric Interpretation}

Salience is interpreted as a curvature-sensitive signal on the semantic
manifold. High novelty corresponds to movement into regions of low prior
density, while high temporal change corresponds to steep gradients in the
trajectory.

This geometric formulation enables downstream components—such as stability
filtering, membrane evaluation, and semantic routing—to respond appropriately
to surprising or rapidly evolving inputs.

Salience acts as the trigger for cognitive updates, ensuring that the system
allocates computational resources to states that are both informative and
behaviorally relevant.
